\documentclass[12pt]{article}
\usepackage{hyperref}
\hypersetup{hidelinks}
\title{LaTeX Drill}
\author{}
\date{}

\begin{document}
\HCode{<a class="skip-link" href="&num;main">本文へスキップ</a>}
\HCode{<header class="site-header"><nav class="site-nav" aria-label="メインナビゲーション"><a class="site-brand" href="https://github.com/Nyanziba/latex-drill">latex-drill</a><div class="nav-links"><a href="&num;curriculum">カリキュラム</a><a href="&num;install">はじめかた</a><a href="&num;sources">公式資料</a></div></nav></header><main id="main" class="site-main">}

\maketitle

\begin{quote}
公式資料を読み、手を動かして身につけるLaTeX練習ドリルです。小さな課題を解いて、コンパイルと課題条件の確認で進み具合を確かめます。
\end{quote}

\HCode{<div class="hero-actions">}
\href{#curriculum}{\textbf{カリキュラムを見る}}
\qquad
\href{https://github.com/Nyanziba/latex-drill}{\textbf{GitHubでドリルを開く}}
\HCode{</div>}

\begin{itemize}
  \item 初級・中級・上級の3トラック
  \item 公式資料に対応した全9題
  \item Windows / Linux / WSLで利用可能
\end{itemize}

\hypertarget{curriculum}{}
\section*{3段階、9つの課題}

課題には読み物、TODO付きのLaTeXファイル、ヒント、解答例があります。各トラックの題材を選び、ドリルのコマンドで読み物と練習を行き来できます。

\HCode{<div class="track-grid"><article class="track-card">}
\subsection*{初級}
文書の骨組みから数式まで、LaTeXの基本を順に練習します。
\begin{itemize}
  \item 文書構造と和文の組版
  \item 論理構造、強調、リスト
  \item 数式モードと基本数式
\end{itemize}
\href{https://github.com/Nyanziba/latex-drill/tree/main/docs/beginner}{初級の読み物を開く}
\HCode{</article><article class="track-card">}

\subsection*{中級}
レポートで使う図表、参照、文献の管理を身につけます。
\begin{itemize}
  \item 図表を読みやすく配置する
  \item ラベルと相互参照を解決する
  \item 引用と参考文献を管理する
\end{itemize}
\href{https://github.com/Nyanziba/latex-drill/tree/main/docs/intermediate}{中級の読み物を開く}
\HCode{</article><article class="track-card">}

\subsection*{上級}
再利用可能なコマンドや小さなパッケージを作ります。expl3を使った状態管理も扱います。
\begin{itemize}
  \item コマンドと環境を定義する
  \item 独自の \texttt{.sty} パッケージを作る
  \item expl3の変数と関数を使う
\end{itemize}
\href{https://github.com/Nyanziba/latex-drill/tree/main/docs/advanced}{上級の読み物を開く}
\HCode{</article></div>}

\hypertarget{install}{}
\section*{はじめかた}

Ubuntu / WSLではリポジトリを取得し、インストーラーを実行します。

\begin{verbatim}
git clone https://github.com/Nyanziba/latex-drill.git
cd latex-drill
./install.sh
./drill list
./drill watch b01
\end{verbatim}

WindowsではPython 3.9以降とTeX Liveを用意して、PowerShellから次を実行します。

\begin{verbatim}
.\drill.bat list
.\drill.bat watch b01
\end{verbatim}

\subsection*{解いて確認する}

課題ファイルの \texttt{\% TODO} を解決して保存すると、自動で条件確認とコンパイルが走ります。すべての条件を満たしたら \texttt{\% I AM NOT DONE} を削除します。まず一覧と対応する読み物を確認してください。

\begin{verbatim}
./drill list
./drill read b01
\end{verbatim}

\hypertarget{sources}{}
\section*{公式資料から学ぶ}

練習内容はLaTeX ProjectやAMSが公開する資料をもとに構成しています。原文の転載ではなく、課題に必要な要点を日本語で説明し、各レッスンから公式の参照先へ移動できるようにしています。

\begin{itemize}
  \item \href{https://www.learnlatex.org/ja/}{Learn LaTeX 日本語版}: 初級の入門レッスン
  \item \href{https://www.latex-project.org/help/documentation/amsldoc_jpn.pdf}{AMS amsmathユーザーガイド}: 数式と中級課題
  \item \href{https://www.latex-project.org/help/documentation/}{LaTeX Project文書一覧}: 著者・クラス・パッケージ開発者向け資料
  \item \href{https://tug.org/tex4ht/}{TeX4ht}: LaTeX文書からHTMLを生成するツール
\end{itemize}

\HCode{</main><footer class="site-footer"><p>サイト原稿はLaTeXで管理し、TeX4htでHTMLに変換しています。</p><a href="https://github.com/Nyanziba/latex-drill">GitHubリポジトリ</a></footer>}
\end{document}
